\documentclass[11pt]{article}
    \usepackage[utf8]{inputenc}
    \usepackage[T1]{fontenc}
    \usepackage[margin=1in]{geometry}
    \usepackage{hyperref}
    \usepackage{enumitem}
    \title{Code.org- Assign Course D Lesson 3 Intro. to Online Puzzles}
    \author{Piter Garcia}
    \date{March 20, 2026}
    \begin{document}
    \maketitle
    \section*{Quick Scan}
    \begin{itemize}[leftmargin=*]
    \item Grade: Grade 4
    \item Workbook sessions: Session 23
    \item Class blocks: Day 1 1:15-2:10 Baris, Day 2 1:15-2:10 Fowler, Day 3 1:15-2:10 Schrank, Day 4 1:15-2:10 Autore
    \item Reference file: code-org-assign-course-d-lesson-3-intro-to-online-puzzles.bib
    \end{itemize}
    \section*{School Planning Goals and Inclusion}
    \begin{itemize}[leftmargin=*]
    \item What is expected: - I can decompose a problem into smaller named parts | - I can develop algorithms or programs that include repetition and conditionals for the purposes of debugging | - I can model how data is structured to transmit through a network
    \item What we achieved: This lesson points to local transcript or evidence files in the workspace so the packet can be revised from what actually happened in class.
    \item What needs improvement: stronger proof, cleaner assessment notes, and tighter lesson artifacts.
    \item How to improve: add transcript proof, one saved artifact, and clearer success criteria.
    \item Inclusion focus: use visual modeling, chunked directions, predictable routines, and flexible participation roles.
    \end{itemize}
    \section*{Official References}
    \begin{itemize}[leftmargin=*]
    \item \url{https://csteachers.org/k12standards/interactive/}
\item \url{https://code.org/curriculum/elementary-school}

    \end{itemize}
    \end{document}
